\chapter{Dissipation from a finite-lag correlation defect}
\label{ch:pa-spectral-dynamics}

A positive departure from a reversible correlation inequality contains more
information than a generic distance-to-equilibrium argument retains. The
operator identity below turns that departure directly into a lower bound on
stationary entropy production. The improvement is linear in a small positive
correlation defect; the previous comparison followed by Pinsker's inequality
was quadratic in its confidence gap. Both statements remain valid under their
respective contracts.

This is a theorem for a specified stochastic model class. It does not identify
the hidden generator, establish the exact minimum dissipation over all
compatible continuous-time models, or infer physical entropy from an arbitrary
oscillatory record. The distinction between ordinary and generalized time
reversal is essential.

\section{Population theorem and its operational assumptions}

Let $X_t$ be a stationary, time-homogeneous Markov process with invariant law
$\pi$ and transition semigroup $P_t$. Initially consider a finite irreducible
CTMC with positive stationary probabilities. Its physical time reversal must
be the \emph{ordinary reversal of the complete hidden state}, so that its
reverse generator is the $L^2(\pi)$ adjoint $Q^\dagger$. Define
\[
 \sigma=\sum_{x\ne y}\pi_x Q_{xy}
       \log\frac{\pi_xQ_{xy}}{\pi_yQ_{yx}}.
\]
The usual conventions give infinite entropy production for an unpaired
positive transition. The finite-entropy case is the informative one below.
Physical heat or power additionally requires a thermodynamic embedding; all
rates here are dimensionless entropy per unit physical time.

Let $f$ be a real bounded function with a valid range bound
$\sup f-\inf f\leq R$, $R>0$. This may be the conditional mean of a
noninvasive memoryless detector. Detector noise must be conditionally
independent between retained readings. Its range bound concerns the whole
admissible channel, not the largest and smallest values in a finite sample.
Set
\[
 C(t)=\langle f,P_t f\rangle_\pi,
 \qquad
 \Delta_s(t)=-s^2 C(0)+2sC(t)-C(2t),\quad t>0.
\]
The population identity allows any real $s$. Selecting $s$ from inference
data requires a simultaneous confidence argument or an independent pilot.

\begin{theorem}[Finite-lag dissipation bound]
\label{thm:pa-finite-lag}
Put $d=[\Delta_s(t)]_+$ and $u=4d/R^2$. Then $0\leq u\leq1$.
If $\sigma$ is finite, $u<1$ whenever $d>0$, and
\begin{equation}
 \boxed{\displaystyle
 \sigma\ \geq\ \frac{\Phi(u)}{t}\ \geq\
 \frac{8[\Delta_s(t)]_+}{R^2 t},
 \qquad
 \Phi(u)=2\sqrt{u}\,\operatorname{atanh}\sqrt{u}.}
 \label{eq:pa-direct-bound}
\end{equation}
Here $\Phi(0)=0$ and $\Phi(1)=+\infty$. The theorem requires no bound on
the number of hidden states or transition rates.
\end{theorem}

\begin{proof}
Write $P=P_t$, $H=(P+P^\dagger)/2$ and $B=(P-P^\dagger)/2$.
These are the self-adjoint and skew-adjoint parts of the transition operator,
not of the generator. In a real Hilbert space,
\begin{align}
 C(2t)&=\langle P^\dagger f,Pf\rangle_\pi
       =\|Hf\|_\pi^2-\|Bf\|_\pi^2,\\
 \Delta_s(t)&=\|Bf\|_\pi^2-\|(H-sI)f\|_\pi^2
              \leq\|Bf\|_\pi^2.
 \label{eq:pa-defect-identity}
\end{align}
In particular a reversible process, for which $B=0$, has no positive defect.

For each state $x$, let $v_x$ be the total variation between the probability
rows $P(x,\cdot)$ and $P^\dagger(x,\cdot)$. The bounded-expectation
inequality gives
\[
 |Bf(x)|\leq\frac{R}{2}v_x,
 \qquad u\leq\sum_x\pi_x v_x^2\leq1.
\]
Let $\mu(x,y)=\pi_xP(x,y)$ and $\mu^\leftarrow(x,y)=\mu(y,x)$ be
the forward and reverse stationary endpoint laws. Transposition exchanges the
two laws, so their relative entropies in opposite directions are equal. Since
both first marginals are $\pi$,
\begin{align}
 D(\mu\Vert\mu^\leftarrow)
 &=\frac12\sum_x\pi_x\bigl\{
       D(P_x\Vert P_x^\dagger)+D(P_x^\dagger\Vert P_x)\bigr\}\\
 &\geq\sum_x\pi_x\,2v_x\operatorname{atanh}v_x.
\end{align}
The last inequality is the established sharp Jeffreys--total-variation
inequality \cite{paSason2015}. The power series
\[
 \Phi(u)=2\sum_{k=1}^\infty\frac{u^k}{2k-1}
\]
shows that $\Phi$ is increasing and convex on $[0,1)$, with
$\Phi(u)\geq2u$. Jensen's inequality therefore gives
\[
 D(\mu\Vert\mu^\leftarrow)
 \geq\Phi\!\left(\sum_x\pi_xv_x^2\right)\geq\Phi(u).
\]
Finally, retaining only the endpoints cannot increase relative entropy.
The forward and ordinarily reversed stationary CTMC path laws over $[0,t]$
have relative entropy $t\sigma$. Consequently
$t\sigma\geq D(\mu\Vert\mu^\leftarrow)\geq\Phi(u)$.
\end{proof}

The elementary linear version follows already from conditional Pinsker and
the same identity. The nonlinear version retains the sharp symmetric
divergence ingredient. Neither version assumes that the observed scalar
record is itself Markovian: the semigroup acts on the full hidden state.

\section{What is sharp, and what remains a realization problem}

The endpoint bound is sharp over stationary discrete-time Markov kernels.
To see this, take a four-state cycle with uniform stationary law and
\[
 P_{i,i+1}=\frac{1+z}{2},\qquad
 P_{i,i-1}=\frac{1-z}{2},\qquad 0<z<1,
\]
with indices modulo four. For $f=(1,1,-1,-1)$ and $s=0$,
\[
 C_0=1,\quad C_1=0,\quad C_2=-z^2,\quad
 \Delta_0=z^2,\quad R=2.
\]
Every row has total variation $z$ from its reverse, and
\[
 D(\mu\Vert\mu^\leftarrow)
 =z\log\frac{1+z}{1-z}=\Phi(z^2).
\]
Thus equality holds throughout the endpoint argument.

This example does \emph{not} establish sharpness among finite CTMCs:
$\det P=-z^2<0$, whereas $\det e^{tQ}=e^{t\operatorname{tr}Q}>0$.
The sharp kernel cannot be the finite-time transition matrix of a finite
CTMC. The contraction from a full continuous path to its endpoints can lose
further information. Determining the best CTMC-compatible rate from a given
finite correlation record remains a separate physical realization problem.

For a nearest-neighbor ring with uniform stationary law and the cosine
detector, write its Fourier eigenvalue as $-\gamma+i\omega$. Then
\[
 C(t)=\tfrac12 e^{-\gamma t}\cos(\omega t),\qquad
 s=e^{-\gamma t}\cos(\omega t),\qquad
 \Delta_s(t)=\tfrac12e^{-2\gamma t}\sin^2(\omega t).
\]
The nonnegative residual in \eqref{eq:pa-defect-identity} vanishes for this
mode, but the subsequent information inequalities need not be saturated.
This avoids interpreting a convenient exact spectral example as a universal
minimum-dissipation model.

\section{Use with finite-confidence calibration}

Suppose an existing confidence construction returns $g\geq0$ such that,
on an event of probability at least $1-\alpha$,
\[
 2g\leq[\Delta_s(t)]_+.
\]
Both the same-state binary-twin construction and the bounded analog joint
construction provide such a gap. For the analog detector, writing
$f(x)=\mathbb E[Z\mid X=x]$ gives
\[
 B_{\rm obs}=\mathbb E Z^2\geq\mathbb E f(X)^2=C(0).
\]
The coefficient of the diagonal moment is $-s^2$, so its excess measurement
noise makes the contrast conservative. With
$q_0=(1+B_{\rm obs})/2$, $q_1=(1+C(t))/2$, $q_2=(1+C(2t))/2$,
\[
 a_s=(-s^2,2s,-1),\qquad b_s=-\frac{(1-s)^2}{2},
 \qquad 2(a_s\cdot q-b_s)\leq\Delta_s(t).
\]
Consequently the same confidence event yields
\begin{equation}
 \sigma\geq\frac{1}{t}\Phi\!\left(\frac{8g}{R^2}\right)
 \geq\frac{16g}{R^2t}.
 \label{eq:pa-statistical-direct}
\end{equation}
For a detector mean in $[-1,1]$, $R=2$ gives the particularly simple
lower bound $4g/t$. There is no new statistical failure allowance for replacing
the old physical implication by this stronger one. Taking a maximum with
another bound is justified without an extra penalty only when both follow
on the same already-covered event, or have separate simultaneous coverage.

The implementation returns exact rational bounds. For $0\leq u<1$ and
an integer $K\geq1$,
\[
 2\sum_{k=1}^K\frac{u^k}{2k-1}
 \leq\Phi(u)\leq
 2\sum_{k=1}^K\frac{u^k}{2k-1}
 +\frac{2u^{K+1}}{(2K+1)(1-u)}.
\]
The upper endpoint encloses the \emph{lower-bound formula}, not the actual
entropy production. If a confidence gap implies $u\geq1$, the API flags
incompatibility with finite entropy \emph{on the coverage event}; it retains
the valid linear inequality and does not silently clip $u$ or treat a possible
confidence failure as an empirical refutation of the model.

\section{Reanalysis of the frozen joint-calibration experiment}

The sixth-edition joint-calibration count archive contains 400 independent
simulated datasets in each of four declared cases. Its contrast, confidence
level, counts, and acquisition resources were fixed before those datasets
were generated. Recomputing the old certificates reproduces all saved rates
exactly. Applying \eqref{eq:pa-statistical-direct} then changes only the
mathematical physical bound, without collecting data or selecting a new
contrast.

\begin{center}
\begin{tabular}{lrrr}
Case & Positive datasets & Previous mean bound & New nonlinear mean bound\\
\hline
Driven ring & $399/400$ & $0.0008005133$ & $0.1122449304$\\
Equilibrium & $0/400$ & $0$ & $0$\\
High readout noise & $6/400$ & $0.0000002836$ & $0.0002243291$\\
Invalid calibration & $400/400$ & $0.0105740727$ & $0.4597513517$\\
\end{tabular}
\end{center}

The driven mean is $140.216$ times the former mean, while its actual model
entropy production is about $3.55736$. The new mean is about $3.16\%$ of
that value: the improved inference still leaves a substantial gap. Positive
detection counts are unchanged because the same positive confidence gap is
used. Each dataset still has 20,000 pair acquisitions, 40,000 endpoint
readouts and 24,000 units of declared evolution time; preparation and detector
overheads require their existing separate calibration.

The invalid-calibration row is a failure control, not a successful inference.
Changing the detector between calibration and science measurements violates
the shared-channel premise and produces false positive certificates even
for equilibrium. A stronger valid inequality cannot repair an invalid
measurement contract. The archived reanalysis is retrospective simulated-data
evidence, not new physical experimental validation.

\section{Vector observations and diffusion processes}

The same proof holds for a bounded detector map $F$ into a real Hilbert space,
with image diameter at most $R$, if
\[
 C(t)=\mathbb E\langle F(X_0),F(X_t)\rangle.
\]
For each row, the difference between the two vector expectations is total
variation times the difference of two barycenters of the detector image.
Its norm is at most $R v_x$, since taking the convex hull does not enlarge
diameter. The operator identity holds componentwise, or directly in the
corresponding Hilbert-valued $L^2$ space. All remaining inequalities are
unchanged. This proves a joint-diameter bound, rather than adding scalar
entropy lower bounds, an addition which would generally be invalid.

For example, replacing the cosine channel by
$F(\theta)=(\cos\theta,\sin\theta)$ doubles the ring correlation and
its defect while retaining diameter two. The linear population bound doubles;
the nonlinear one improves further. This is a changed observation channel,
whose full vector acquisition and readout costs must be recorded. It is not
additional information obtained by processing the original cosine alone.
For a noisy vector detector $Z$ confined to the unit ball,
$\mathbb E\|Z\|^2\geq\mathbb E\|\mathbb E(Z\mid X)\|^2$, and the
scores $(1+\|Z\|^2)/2$ and $(1+\langle Z_0,Z_t\rangle)/2$ have ranges
$[1/2,1]$ and $[0,1]$. This supplies the analogous joint-score construction;
the executable scalar benchmark does not claim a measured vector advantage.

The result also transfers to a stationary ordinary-even overdamped diffusion
whenever the forward and reversed path laws exist with
$D(\mathbb P_{[0,t]}\Vert\mathbb P^\leftarrow_{[0,t]})=t\sigma$.
For the smooth, compact, uniformly elliptic setting, or the separately
justified absolute-continuity setting of the diffusion chapter, the proof
uses the same transition-kernel identity and conditional divergences.
Here the reverse comparison is $Q^\dagger$, and its path divergence is
$t\sigma$, rather than the earlier reversible comparison's $t\sigma/4$.
Those constants belong to different comparisons and cannot be interchanged.

\section{Counterexamples that delimit the inference}

Stationarity cannot be replaced by a reversible generator alone. Take
$Q=\left(\begin{smallmatrix}-1&1\\1&-1\end{smallmatrix}\right)$,
$f=(1/2,1)$ and initialize the process in state zero. Its nonstationary
products obey $C(t)=3/8-e^{-2t}/8$, so
\[
 -C(0)+2C(t)-C(2t)=\tfrac18(1-e^{-2t})^2>0.
\]
The generator's stationary entropy production is zero. The positive record
defect therefore does not imply a positive stationary entropy rate without
the stationary-preparation assumption.

Likewise, observing an even physical quantity is insufficient to establish
ordinary time reversal for the hidden state. An equilibrium inertial
oscillator has an even position and odd velocity; its bounded position
readouts can have oscillatory correlations and a positive defect while its
physical entropy production is zero. Another explicit model is circle
drift--diffusion with the physical involution $\theta\mapsto-\theta$.
Its generator $Q=\omega\partial_\theta+D\partial_\theta^2$ obeys
$\Theta Q\Theta=Q^\dagger$ and is generalized reversible under that
involution. The even observable $\cos\theta$ has
$C(t)=e^{-Dt}\cos(\omega t)/2$; at $\omega t=\pi/2$ and $s=0$,
$\Delta_s=e^{-2Dt}/2>0$. Its ordinary-reversal mathematical entropy and
its generalized-reversal physical entropy are different quantities. The
theorem applies to the former unless ordinary full-state reversal has been
physically justified.

Finally, related spectral, cross-correlation and variance-sum-rule
dissipation bounds are established prior work. The specific contribution
tested here is the finite-lag scalar-defect implication, its sharp endpoint
envelope with explicit nonembedding caveat, and its integration with the
already-covered calibration record. The adjacent literature assessment
identifies the closest results and does not claim historical priority.
